\section{An upper bound from a signed variational expression}
\label{S:original-signed-score-upper}

Keep all auxiliary choices fixed as in the Gaussian tube construction, including
$\alpha/4 < s_{\rm sp} < 1/2$, the admissible $\theta,\nu,r,p$, and the three
fixed even-$Q$ conditions. The mean bound $C_{\rm mean}$ is independent of
$R,B,\delta$; it may depend on those fixed auxiliary parameters and the model.
The space $H_I$ has the original covariance $\Gamma$. Its norm is not reweighted
by $g$ and is not a sum of fast block costs. Let $R_\delta = \delta^{-1}+2$ and
$M_\delta = \E\|P_\delta\|_{\mathcal{L}(E)}^2$. For every fixed integer $B$ the
Gaussian block tuple has typical radius $b_{\delta,B} =
2\E\|X_{\delta,B}\|_{\delta,p} \to 0$. Its endpoint is the same $P_\delta$ for
every $B$ by Lemma~\ref{L:tube-actual-endpoint}, not a $B$-dependent surrogate.
Write $\mathcal{M}_{\delta,B,R}$ for the normalized mean tuple
in~\eqref{E:tube-actual-mean}, with the indicated block length and radius. Write
$\mathcal{F}_{\delta,B}$ for the everywhere-defined endpoint functional
$F_\delta$ of Lemma~\ref{L:tube-actual-endpoint}.

\begin{lemma}[Signed variational expression and quantitative radial envelope]
  \label{L:signed-score-envelope}
  For $R \ge 1$ define the variational expression
  \begin{equation}\label{E:signed-raw-score}
    \mathcal{S}_R=\sup_{\Phi\in H_I}
    \{2\operatorname{Rate}_R(\Phi)-\|\Phi\|_{H_I}^2/2\},
    \qquad V_{\rm up}=\limsup_{\delta\downarrow0}\mathcal{S}_{R_\delta}.
  \end{equation}
  Then $0 \le \mathcal{S}_R,V_{\rm up} \le 2+2C_{\rm mean}^2 < \infty$. The tube
  envelope $G_{\delta,B}$ of Proposition~\ref{P:tube-coarse-envelope} satisfies,
  uniformly in $R \ge 1$ (in particular at the stochastic radius $R_\delta$),
  \begin{equation}\label{E:signed-polynomial-envelope}
    G_{\delta,B}(\rho)\le A_0+A_1B^{3\nu}\rho^\nu
    +\eta_B b_{\delta,B}^\nu ,
  \end{equation}
  for $0 < \delta < \min(1,1/B)$ and every $\rho \ge 0$, where
  \[
    \begin{aligned}
      A_0&=4C_{\rm gr},\\
      A_1&=2^{\nu+1-\nu/r}C_{\rm gr}32^\nu C_{\rm mean}^\nu,\\
      \eta_B&=2^{\nu+1-\nu/p}C_{\rm gr}B^{\nu/p}.
    \end{aligned}
  \]
  The variational expression is not identified here with a constant in the cited
  source or a variational value over nonnegative controls.
\end{lemma}

\begin{proof}
  We bound the signed score by completing the square in the control norm. We
  then use the explicit physical Young growth bound and weighted block
  estimates to compute the radial envelope constants, keeping their dependence
  on $B$.


  \medskip
  \noindent\textit{Step 1.~}We obtain a finite uniform ceiling for the signed
  score.
  The mean is linear, so $\mathcal{T}_R0 = 0$. Its free form has nonnegative
  gradient energy and spectral bottom zero; thus $\operatorname{Rate}_R(0) = 0$.
  Proposition~\ref{P:uniform-rate-averaging} gives $\operatorname{Rate}_R(\Phi)
  \le 1+C_{\rm mean}\|\Phi\|$. With $q = \|\Phi\| \ge 0$,
  \[
    2\operatorname{Rate}_R(\Phi)-q^2/2
    \le2+2C_{\rm mean}q-q^2/2
    =2+2C_{\rm mean}^2-(q-2C_{\rm mean})^2/2 .
  \]
  The coarser inequality also follows from half-weight arithmetic--geometric
  mean at bases $q^2,(2C_{\rm mean})^2$, with nonnegative weights and bases,
  including zeros. The zero control gives the lower bound. No optimizer, finite
  rank, positive control or convergence of $\mathcal{S}_{R_\delta}$ is assumed.
  The finite ceiling is auxiliary-parameter dependent and need not be sharp.


  \medskip
  \noindent\textit{Step 2.~}We compute the block-length dependence of the tube
  envelope.
  For $\|a\|_{\delta,p} \le b_{\delta,B}$, $\|\Phi\| \le \rho$ set $v_m =
  \|a_m\|_\theta$ and $h_m = \|(\mathcal{M}_{\delta,B,R}\Phi)_m\|_\theta$. The
  quantitative physical bound~\eqref{E:uniform-product-growth}, already
  including the free factor, and convexity give
  \[
    \mathcal{F}_{\delta,B}(a+\mathcal{M}_{\delta,B,R}\Phi)
    \le2C_{\rm gr}B\delta n
    +2^\nu C_{\rm gr}B\delta\sum_m(v_m^\nu+h_m^\nu).
  \]
  Its upper expression is nonnegative, so it controls the maximum at one. The
  true primitive estimate and weighted power mean give
  \begin{align*}
    \delta\sum v_m^\nu&\le(2/B)^{1-\nu/p}b_{\delta,B}^\nu,\\
    \delta\sum h_m^\nu&\le(2/B)^{1-\nu/r}
    (32B^{3-1/r}C_{\rm mean}\rho)^\nu .
  \end{align*}
  The exponents are $p/\nu \ge 1$, $r/\nu > 1$, inputs nonnegative, $n = N+1 \ge
  2$, $\delta > 0$, and $\delta n \le 2/B$. Zeros and $p = \nu$ are legal. The
  first constant is $4C_{\rm gr}$. In the mean term the factors $B^{\nu/r}$ and
  $B^{-\nu/r}$ cancel exactly, leaving $B^{3\nu}$. These computations give the
  displayed coefficients and~\eqref{E:signed-polynomial-envelope}, uniformly for
  all radii. They use no unspecified $c_Y(B)$ or quantitative $F$ estimate. This
  proves Lemma~\ref{L:signed-score-envelope}.
\end{proof}

Lemma~\ref{L:signed-score-envelope} bounds the envelope $G_{\delta,B}$ at every
radius $\rho$. We now split the radial integral
in~\eqref{E:tube-coarse-envelope} at a radius $M$. Below $M$ we compare the
exponent with $\mathcal{S}_{R_\delta}$, and above $M$ we use
Lemma~\ref{L:tube-rayleigh-tail}.

\begin{proposition}[Finite-delta bulk and exact tail]
  \label{P:signed-moment-bulk-tail}
  For fixed integer $B \ge 1$, $0 < \delta < \min(1,1/B)$ and $M > 0$, define
  the deterministic number
  \[
    \mathcal{E}_{\delta,B,M}=
    \max\left(0,\sup_{\substack{R\ge1,\ \|\Phi\|\le M\\
    \|x\|_{\delta,p}\le b_{\delta,B}}}
    [\mathcal{F}_{\delta,B}(x+\mathcal{M}_{\delta,B,R}\Phi)
    -2\operatorname{Rate}_R(\Phi)]\right).
  \]
  It is finite. For every polynomial $M_B = c_{\rm ball}B^{\kappa_{\rm ball}}$,
  with $c_{\rm ball} > 0$, $\kappa_{\rm ball} \ge 0$, one has
  \begin{equation}\label{E:signed-error-limit}
    e_B:=\limsup_{\delta\downarrow0}\mathcal{E}_{\delta,B,M_B}
    \ \text{finite and nonnegative},\qquad e_B\longrightarrow0 .
  \end{equation}
  If $4A_1B^{3\nu} \le M^{2-\nu}$, then
  \begin{equation}\label{E:signed-finite-moment}
    M_\delta\le
    \frac{M^2}{2\delta}
    e^{(\mathcal{S}_{R_\delta}+\mathcal{E}_{\delta,B,M})/\delta}
    +4e^{(A_0+\eta_Bb_{\delta,B}^\nu-M^2/4)/\delta}.
  \end{equation}
\end{proposition}

\begin{proof}
  We first identify the vanishing deterministic error on polynomial control
  balls. We then split the exact radial integral at $M$: the bulk uses the
  signed variational score, and the tail uses the explicit subquadratic
  envelope.


  \medskip
  \noindent\textit{Step 1.~}We control the deterministic product error.
  The polynomial envelope and nonnegative rate make this supremum finite. It is
  an ordinary deterministic real number, not a random uncountable-radius
  supremum. At fixed $R$ the inner functions are continuous in controls and
  paths; no measurable maximizing selector or joint $\delta,R$ measurability is
  needed. The zero control and zero path have rate zero and
  $\mathcal{F}_{\delta,B}(0) \ge 0$, so the untruncated supremum is already
  nonnegative. The anchor $b_{\delta,B} \to 0$ at fixed $B$ satisfies
  Proposition~\ref{P:uniform-arbitrary-residual} exactly. Its polynomial-radius
  conclusion gives $\limsup_B e_B \le 0$. At fixed $B$, the polynomial envelope
  bounds the inner limsup by a finite number, and $e_B \ge 0$ then
  proves~\eqref{E:signed-error-limit}.


  \medskip
  \noindent\textit{Step 2.~}We split the radial integral and keep both
  prefactors.
  The exact Gaussian-enlargement bound is
  \[
    M_\delta\le\int_0^\infty\frac{\rho}{\delta}
    e^{(G_{\delta,B}(\rho)-\rho^2/2)/\delta}\ud\rho .
  \]
  For every $\rho \le M$, every base-ball $a$ and $\|\Phi\| \le \rho$,
  \[
    \mathcal{F}_{\delta,B}(a+\mathcal{M}_{\delta,B,R_\delta}\Phi)-\rho^2/2
    \le2\operatorname{Rate}_{R_\delta}(\Phi)-\|\Phi\|^2/2
    +\mathcal{E}_{\delta,B,M}
    \le\mathcal{S}_{R_\delta}+\mathcal{E}_{\delta,B,M}.
  \]

  Taking suprema and integrating $\rho/\delta$ over $[0,M]$ gives the first term
  of~\eqref{E:signed-finite-moment}, even if $M^2/(2\delta) < 1$. This used
  $\rho^2 \ge \|\Phi\|^2$, not equality or an optimizer. For the tail apply
  Lemma~\ref{L:tube-rayleigh-tail} with $C = A_1B^{3\nu}$ and intercept
  $A_0+\eta_Bb_{\delta,B}^\nu$. All cutoff conditions hold: $M > 0$, $\rho \ge
  M$, $C \ge 0$, $1 < \nu < 2$, $4C \le M^{2-\nu}$. Its analytic integration
  retains both the front factor four and the exponent $-M^2/4$. Adding bulk and
  tail proves the finite-delta bound. No infinite-dimensional density or Laplace
  principle is used. This proves Proposition~\ref{P:signed-moment-bulk-tail}.
\end{proof}

\begin{proposition}[Two limits on the same moment]
  \label{P:signed-operator-upper}
  The global operator moment obeys
  \[
    0\le\liminf_{\delta\downarrow0}\delta\log M_\delta
    \le\limsup_{\delta\downarrow0}\delta\log M_\delta
    \le V_{\rm up}\le2+2C_{\rm mean}^2 .
  \]
\end{proposition}

\begin{proof}
  We choose a polynomial cutoff large enough for the quadratic tail to
  dominate. We first take the small-parameter limit at fixed block length, and
  only then send the block length to infinity, keeping the same endpoint
  moment on the left.


  \medskip
  \noindent\textit{Step 1.~}We choose a polynomial tail cutoff and bound the
  logarithm of the sum.
  Choose any $c_{\rm ball} > 0$ and $\kappa_{\rm ball} > 3\nu/(2-\nu)$, and set
  $M_B = c_{\rm ball}B^{\kappa_{\rm ball}}$. The cutoff in
  Proposition~\ref{P:signed-moment-bulk-tail} holds for all sufficiently large
  integers $B$. For $A_1 = 0$ this is immediate; otherwise its side ratio is
  \[
    \frac{M_B^{2-\nu}}{4A_1B^{3\nu}}
    =\frac{c_{\rm ball}^{2-\nu}}{4A_1}
    B^{\kappa_{\rm ball}(2-\nu)-3\nu}\longrightarrow\infty .
  \]
  Both constants and the strictly positive exponent are retained. At such a
  fixed $B$, a positive sum is at most twice its larger term. Taking
  $\delta\log$ of~\eqref{E:signed-finite-moment} therefore gives
  \begin{align*}
    \delta\log M_\delta\le\delta\log2+
    \max\{&
    \mathcal{S}_{R_\delta}+\mathcal{E}_{\delta,B,M_B}
    +\delta\log(M_B^2/(2\delta)),\\
    &A_0+\eta_Bb_{\delta,B}^\nu-M_B^2/4+\delta\log4\}.
  \end{align*}

  \medskip
  \noindent\textit{Step 2.~}We take the small-parameter limit before the
  block-length limit.
  At fixed $B$, $b_{\delta,B} \to 0$. All logarithmic normalizers vanish,
  including $\delta\log(M_B^2/(2\delta)) = \delta\log(M_B^2/2)-\delta\log\delta
  \to 0$. The last limit follows from $te^{-t} \le 2/t$ for $t = -\log\delta >
  0$. Finite variational expression and error bounds permit limsup
  subadditivity, giving
  \[
    \limsup_{\delta\downarrow0}\delta\log M_\delta
    \le\max\{V_{\rm up}+e_B,\ A_0-M_B^2/4\}.
  \]

  The lower bound $M_\delta\ge1$ used here is~\eqref{E:point-fixed-finiteness}
  of Lemma~\ref{L:point-actual-finiteness}.
  The left side has no $B$: restart identifies every block product with this
  same $P_\delta$. Now let integer $B \to \infty$. We have $e_B \to 0$, $M_B \to
  \infty$ and finite $V_{\rm up} \ge 0$. This gives the asserted upper bound,
  while $M_\delta \ge 1$ gives the nonnegative lower bound. The order was
  $\delta$ first at fixed block length, ball and thresholds, then $B$; no
  diagonal or limit of $\mathcal{S}_{R_\delta}$ was assumed. This proves
  Proposition~\ref{P:signed-operator-upper}.
\end{proof}

\begin{proposition}[Pointwise upper bound from the signed variational expression]
  \label{P:original-signed-point-upper}
  For every original admissible model and every fixed $x$,
  \[
    \limsup_{T\to\infty}T^{-\kappa}
    \log\E|u_\Gamma(T,x)|^2
    \le\limsup_{\delta\downarrow0}\mathcal{S}_{\delta^{-1}+2}
    \le2+2C_{\rm mean}^2 .
  \]
  For zero covariance the two normalized moments and the signed variational
  expression have rate zero.
\end{proposition}

\begin{proof}
  Proposition~\ref{P:original-point-transfer} bounds the original point moment
  by $\operatorname{DataPref}_\delta M_\delta$, with
  $\delta\log\operatorname{DataPref}_\delta \to 0$. Its with the regularization
  removed first, full-chain and canonical law proof is the input here, not an
  original $E$ point trace. The decreasing bijection $T_\delta =
  \delta^{-1/\kappa}$ has inverse $\delta = T^{-\kappa}$. Take logs, normalized
  limsup and Proposition~\ref{P:signed-operator-upper}. Original canonical
  spatial stationarity from Proposition~\ref{P:raw-block-original-lower} gives
  every fixed $x$. The moment is at least one by its full expansion and finite
  by the bound, so no logarithm or difference of infinite quantities is hidden.

  For $\Gamma = 0$ set $C_{\rm mean} = 0$ directly. Then $H_I = \{0\}$,
  $\operatorname{Rate}_R = \mathcal{S}_R = V_{\rm up} = 0$, $u = 1$ and
  $P_\delta = S_E(\delta^{-1})$. The free lower and upper bounds give $0 \le
  \delta\log M_\delta \le 2\delta\log[C(1+\delta^{-1})] \to 0$. This direct
  proof needs no nonzero Gaussian or nondegenerate Cameron assumption; the free
  functional at zero may exceed zero. The formal $a\downarrow0$ exponent is
  $(4-\alpha)/(3-\alpha)$, with no new endpoint-in-$a$ theorem asserted. The
  Jensen and lower bounds remain valid. This proves
  Proposition~\ref{P:original-signed-point-upper}.
\end{proof}

The varying variational expression remains a signed variational value for the
original covariance. Its positive-control comparison, exact limit and
identification of the derived constant require the subsequent variational
arguments.
